\section{Conclusion}

The main result is conditional but exact. The paper does not claim that quantum mechanics requires
new dynamics, nor that Everettian physics is false. It claims something narrower and more basic.
If a physically non-selective theory is used to underwrite truth-apt diachronic discourse in which
measurement records function as evidential anchors, then admissible reference assignments for those
records must collapse into a single standing-bearing class on the fixed scope. If they do not, then
confirmation, memory, anticipation, and deliberation become assignment-relative rather than
 theory-relative.

The theorem ladder reaches that conclusion through four steps. First, admissible reference assignments
are typed in fixed-domain form and their standing-bearing content is shown to factor through a unique
measurement-reference quotient. Second, the bare Everettian package is shown to be physically
non-selective relative to the target discourse. Third, theory-level confirmation is shown to require
fixation commutation under admissible redescription; quotient-relevant divergence therefore makes
confirmation undefined on the fixed scope rather than merely weaker. Fourth, memory, anticipation,
deliberation, and confirmation are shown to share one identity-preserving referential lineage through
the relevant measurement records. The interpretation taxonomy is then a downstream exhaustion
result about where that lineage can be fixed, if it is fixed at all.

The resulting diagnosis differs from the standard presentation of the measurement problem. On the
target discourse, the first failure is not yet that a physically non-selective theory does not produce one
outcome. The first failure is that the theory may leave underfixed the referents of the very records
through which outcomes are reported, remembered, anticipated, and used as evidence. In that precise
sense, the measurement problem is a problem of reference to measurement records before it is a
problem of outcome production.
